\slajdid{02-s034}
\begin{frame}[label=02-s034]{Karnoove karte jedne i dve promenljive}
\setlength{\parskip}{0pt}
% element: 02-s034:e01
\alert{Susedna polja imaju Hemingovo rastojanje 1.}
\par\vspace{3mm}
\begin{columns}[T,onlytextwidth]
\begin{column}{.50\textwidth}\centering
% element: 02-s034:e02
\Oznaka{Tabela sa jednom promenljivom}\vspace{2mm}
\begin{tikzpicture}[font=\fontsize{9}{10.5}\selectfont,x=1mm,y=1mm,
 DEindeks/.style={font=\fontsize{7}{8}\selectfont,anchor=south east,inner sep=.7pt}]
\draw (0,0) rectangle (16,8) (8,0)--(8,8);
\draw (0,8)--(-3,11) node[above] (a) {A};
\node[above] at (4,8) (v0) {0};\node[above] at (12,8) (v1) {1};
\node[DEindeks] at (8,0) (i0) {0};
\node[DEindeks] at (16,0) (i1) {1};
\node at (5,20) (prom) {Promenljiva};
\draw[DEcrvena,-{Latex[length=1.4mm]}] (prom.south)--(a.north);
\node[align=left,anchor=west] at (22,14) (val) {Vrednosti\\promenljive};
\draw[DEcrvena,-{Latex[length=1.4mm]}] (val.west)--(v1.north);
\node at (8,-6) (ind) {Indeksi polja};
\draw[DEcrvena,-{Latex[length=1.4mm]}] (ind.north)--(i0.south);
\draw[DEcrvena,-{Latex[length=1.4mm]}] (ind.north)--(i1.south);
\node at (44,4) {ili};
\draw (54,-4) rectangle (62,12) (54,4)--(62,4);
\draw (54,12)--(51,15) node[above] {A};
\node[left] at (54,8) {0};\node[left] at (54,0) {1};
\node[DEindeks] at (62,4) {0};
\node[DEindeks] at (62,-4) {1};
\end{tikzpicture}

\par\vspace{2mm}
\raggedright
Indeks polja je binarna vrednost stanja promenljivih.
\end{column}
\begin{column}{.46\textwidth}\centering
\Oznaka{Tabela sa dve promenljive}\vspace{2mm}
\begin{tikzpicture}[font=\fontsize{9}{10.5}\selectfont,x=4.8mm,y=4.8mm,
 DEindeks/.style={font=\fontsize{7}{8}\selectfont,anchor=south east,inner sep=.7pt}]
\draw (0,2) rectangle (2,4) (1,2)--(1,4) (0,3)--(2,3);
\draw (0,4)--(-.5,4.5) node[above right] {A} node[below left] {B};
\foreach \x/\lab in {.5/0,1.5/1}{\node[above] at (\x,4) {\lab};}
\foreach \y/\lab in {3.5/0,2.5/1}{\node[left] at (0,\y) {\lab};}
\foreach \x/\y/\lab in {1/3/0,2/3/1,1/2/2,2/2/3}{\node[DEindeks] at (\x,\y) {\lab};}
\node at (1,1) {1};
\node at (2.9,3.2) {ili};
\draw (5,0) rectangle (6,4);
\foreach \y in {1,2,3}{\draw (5,\y)--(6,\y);}
\draw (5,4)--(4.5,4.5) node[above left] {BA};
\foreach \y/\lab in {3.5/00,2.5/01,1.5/11,.5/10}{\node[left] at (5,\y) {\lab};}
\foreach \y/\lab in {3/0,2/1,1/3,0/2}{\node[DEindeks] at (6,\y) {\lab};}
\draw[DEcrvena] (5,0)--(6,0) (5,4)--(6,4);
\node at (6.7,.4) {2};
\node at (7,3.2) {ili};
\draw (8.2,3) rectangle (12.2,4);
\foreach \x in {9.2,10.2,11.2}{\draw (\x,3)--(\x,4);}
\draw (8.2,4)--(7.7,4.5) node[above right] {BA};
\foreach \x/\lab in {8.7/00,9.7/01,10.7/11,11.7/10}{\node[above] at (\x,4) {\lab};}
\foreach \x/\lab in {9.2/0,10.2/1,11.2/3,12.2/2}{\node[DEindeks] at (\x,3) {\lab};}
\draw[DEcrvena] (8.2,3)--(8.2,4) (12.2,3)--(12.2,4);
\node at (10.2,2) {3};
\end{tikzpicture}

\par\vspace{2mm}
\raggedright
\alert{Susedna polja imaju najmanje jednu zajedničku stranicu.}
\par\vspace{2mm}
Naspramne ivice karte se povezuju.
\end{column}
\end{columns}
\end{frame}
